% !TEX root = ../main.tex
\section{Ejercicio propuesto 1: Biblioteca integrada y TCR}

\subsection{Especificación de la biblioteca \texttt{number\_theory.py}}
La biblioteca reúne las funciones de la interfaz requerida en un módulo estructurado con validación estricta de tipos y excepciones, sintetizadas en el Cuadro \ref{tab:resumen_funciones}.

\begin{table}[H]
    \caption{Especificación de la interfaz en \texttt{number\_theory.py}}
    \label{tab:resumen_funciones}
    \centering
    \footnotesize
    \begin{tabular*}{\textwidth}{@{\extracolsep{\fill}}llll}
        \toprule
        \textbf{Función} & \textbf{Retorno} & \textbf{Excepción (\texttt{ValueError})} & \textbf{Complejidad} \\
        \midrule
        \texttt{gcd\_trace(a, b)} & $(d, \text{traza})$ & — & $O(k)$ divisiones \\
        \texttt{extended\_gcd(a, b)} & $(d, x, y)$ con $ax+by=d$ & — & $O(k)$ iteraciones \\
        \texttt{modular\_inverse(a, n)} & $u \in \{1, \dots, n-1\}$ & $n \le 1$; $\gcd(a, n) \neq 1$ & $O(k)$ euclídeo \\
        \texttt{solve\_linear\_congruence} & Lista de $d$ soluciones & $n \le 0$ & $O(k + d)$ \\
        \texttt{modular\_pow(b, e, m)} & $(R, \text{sq}, \text{prod})$ & $m \le 0$; $e < 0$ & $\le 2k - 1$ mult. \\
        \texttt{crt(residues, moduli)} & $(x, N)$ con $0 \le x < N$ & Longitudes, $n_i \le 1$, $\gcd(n_i, n_j) \neq 1$ & $O(m^2 \log n_{\text{máx}})$ \\
        \texttt{validate\_crt\_solution} & Booleano (\texttt{True/False}) & Longitudes distintas; $n_i \le 0$ & $O(m)$ divisiones \\
        \bottomrule
    \end{tabular*}
\end{table}

\subsection{Implementación analítica del Teorema Chino del Resto}
La función \texttt{crt\_with\_trace} resuelve sistemas simultáneos de forma puramente constructiva sin recurrir a búsquedas exhaustivas:

\begin{lstlisting}[language=Python, caption={Teorema Chino del Resto con traza}, label={lst:crt_code}]
def crt_with_trace(residues, moduli):
    residues, moduli = _validate_crt_input(residues, moduli)
    N = prod(moduli)
    total, trace = 0, []
    for i, (a_i, n_i) in enumerate(zip(residues, moduli), start=1):
        a_norm = a_i % n_i
        N_i = N // n_i
        M_i = modular_inverse(N_i % n_i, n_i)
        term = a_norm * N_i * M_i
        total += term
        trace.append({"i": i, "a_i": a_norm, "n_i": n_i, "N_i": N_i, "M_i": M_i, "a_i*N_i*M_i": term})
    return total % N, N, trace
\end{lstlisting}

\subsection{Resolución del caso obligatorio del TCR}
Se evaluó el sistema canónico de la guía: $x \equiv 2 \pmod 3$, $x \equiv 3 \pmod 5$, $x \equiv 2 \pmod 7$. El módulo total es $N = 3 \times 5 \times 7 = 105$. La traza constructiva generada se detalla en el Cuadro \ref{tab:tcr_caso_obligatorio}.

\begin{table}[H]
    \caption{Traza de construcción analítica para el caso obligatorio del TCR}
    \label{tab:tcr_caso_obligatorio}
    \centering
    \footnotesize
    \begin{tabular*}{\textwidth}{@{\extracolsep{\fill}}ccccccc}
        \toprule
        $i$ & Residuo $a_i$ & Módulo $n_i$ & Parcial $N_i = N/n_i$ & Inverso $M_i = N_i^{-1} \bmod n_i$ & Base $e_i = N_i M_i$ & Término $a_i N_i M_i$ \\
        \midrule
        1 & 2 & 3 & 35 & $35^{-1} \bmod 3 = 2$ & 70 & $2 \times 35 \times 2 = 140$ \\
        2 & 3 & 5 & 21 & $21^{-1} \bmod 5 = 1$ & 21 & $3 \times 21 \times 1 = 63$ \\
        3 & 2 & 7 & 15 & $15^{-1} \bmod 7 = 1$ & 15 & $2 \times 15 \times 1 = 30$ \\
        \midrule
        \multicolumn{7}{l}{\textbf{Suma acumulada:} $\sum a_i N_i M_i = 140 + 63 + 30 = 233 \implies x = 233 \bmod 105 = 23$.} \\
        \multicolumn{7}{l}{\textbf{Salida canónica:} $(x, N) = (23, 105)$. \textbf{Comprobación:} $23 \equiv 2 \pmod 3$, $23 \equiv 3 \pmod 5$, $23 \equiv 2 \pmod 7$.} \\
        \bottomrule
    \end{tabular*}
\end{table}

\subsection{Comprobación de propiedades formales en la suite de pruebas}
El archivo \texttt{test\_number\_theory.py} implementa 34 pruebas automatizadas que cubren exhaustivamente las 10 propiedades algebraicas exigidas por la guía (traza euclidiana $a_i = b_i q_i + r_i$, coincidencia de $\gcd$ con \texttt{math.gcd}, identidad $ax+by=d$, inversa válida y rechazo cuando $\gcd \neq 1$, $d$ soluciones canónicas en congruencias, potencia rápida contra \texttt{pow}, unicidad del TCR y normalización en $[0, n-1]$). Todas las pruebas completan con 0 fallos en $\approx 3.4$ segundos.
